\documentclass[11pt,a4paper]{article}
\usepackage[margin=2.6cm]{geometry}
\usepackage{amsmath,amssymb,amsthm}
\usepackage{booktabs}
\usepackage{microtype}
\usepackage{xcolor}
\usepackage[colorlinks=true,linkcolor=blue!50!black,citecolor=blue!50!black,urlcolor=blue!50!black]{hyperref}

\emergencystretch=2em

\newtheorem{theorem}{Theorem}
\newtheorem{lemma}[theorem]{Lemma}
\theoremstyle{remark}
\newtheorem{remark}[theorem]{Remark}

\newcommand{\Lam}{\Lambda}
\newcommand{\eps}{\varepsilon}
\DeclareMathOperator{\Pf}{Pf}

\title{The window quadratic form of the TFPT prime-front program\\
is a Galerkin discretization of Suzuki's localized Weil operator}
\author{Stefan Hamann\thanks{Working note of the verification program
of Topological Fixed-Point Theory (TFPT),
\url{https://fixpoint-theory.com}.  Every number quoted in this note is
produced by a machine-checked script named in the text (repository
file \texttt{verification/v563\_paper2\_readouts.py}, the suite modules
of Remark~\ref{rem:suite}, and the discovery probes
\texttt{experiments/tfpt-discovery/w1\_*.py},
\texttt{w2\_form\_density\_probe.py}); no number is quoted that was not
printed by an executed run.}}
\date{August 2, 2026}

\begin{document}
\maketitle

\begin{abstract}
Suzuki (arXiv:2606.09096) realizes the Weil quadratic form localized to
$[-a,a]$ as $Q_W^a(v)=\langle B_a v,v\rangle$ with
$B_a=D^*G_aD$ on $H^1_0(-a,a)$, where $G_a$ is the integral operator of
the screw function $g$ associated with $\zeta(s)$, and identifies the
Connes--Consani--Moscovici operator $A_a$ as the Friedrichs extension of
$B_a$.  Independently, the prime-front program of Topological
Fixed-Point Theory (TFPT) --- a machine-checked verification suite in
which every quoted number is printed by a named script --- measures
positivity of a finite ``window'' matrix $A_h$ assembled from tent
(hat) reads of the Weil explicit-formula measure on a declared
lattice.  We prove, at the
level of measures and with every layer either symbolic or certified to
$20$--$50$ digits, that these are the same object: the window matrix is
a hat-function Galerkin discretization of $B_a$, with a
\emph{one-scalar} dictionary $c_{\mathrm{TFPT}}(d)=D^{-1}
\bigl(c_{\mathrm{gal}}(d)-\mathrm{pole}_K(d)\bigr)$ up to a derived,
window-independent second-order law $R(d)\to 1+1/(6d^2)$ and an exact
finite-rank pole block.  The $L^2_0$ projection in Suzuki's construction
imposes no condition on the Galerkin coefficients (the mean-zero
condition lives on the differentiated objects and is automatic on
$H^1_0$), and the TFPT parity projection selects an odd subspace of the
form domain, so the bookkeeping closes trivially.  On the declared
window ($h=184$) the two quadratic forms agree on ten odd test vectors
to relative deviation $\le 1.3\times10^{-10}$; the dictionary,
frozen, transports to fresh windows $h=285,540,997$ with only the
explicit lattice constant $D$ changing.  No positivity claim and no
statement about the Riemann Hypothesis is made: the implication chain
beyond the identification (form density, spectral limits, positivity
transfer) is stated as open.
\end{abstract}

\medskip
\noindent\textbf{MSC 2020.} 11M26 (primary); 47B25, 65N30 (secondary).\\
\textbf{Keywords.} Weil explicit formula; screw function; localized
Weil operator; Galerkin discretization; Riemann zeta function;
Hurwitz--Lerch zeta function.

\section{The two objects}

\subsection{Suzuki's localized Weil operator}
Let $\Lam$ be the von Mangoldt function,
$\psi$ the digamma function, $\Phi(z,s,a)=\sum_{n\ge0}(n+a)^{-s}z^n$ the
Hurwitz--Lerch zeta function, and let
\begin{multline}\label{eq:screw}
g(t)=-4\bigl(e^{t/2}+e^{-t/2}-2\bigr)
+\!\!\sum_{n\le e^{|t|}}\!\frac{\Lam(n)}{\sqrt n}\,(|t|-\log n)\\
-\frac{|t|}{2}\Bigl(\psi(\tfrac14)-\log\pi\Bigr)
-\frac14\Bigl(\Phi(1,2,\tfrac14)-e^{-|t|/2}\,\Phi(e^{-2|t|},2,\tfrac14)\Bigr)
\end{multline}
be the screw function associated with $\zeta(s)$
\cite[eq.~(1.3)]{Suz26}.  The final bracket carries the prefactor
$\tfrac14$ on \emph{both} terms; this reading is locked by the paper's
own data (see Remark~\ref{rem:convention}).  In particular $g(0)=0$ and,
near the origin \cite[Sec.~2.2]{Suz26},
\begin{equation}\label{eq:origin}
g(t)=\tfrac12|t|\log|t|+A|t|
+\!\!\sum_{n\le e^{|t|}}\!\frac{\Lam(n)}{\sqrt n}(|t|-\log n)+O(t^2),
\qquad A=\tfrac12\bigl(\log 2\pi-\psi(2)\bigr)=0.7075\ldots
\end{equation}
With $(Gu)(x)=\int g(x-y)u(y)\,dy$, $L^2_0(-a,a)=\{u:\int u=0\}$, $P_a$
the orthogonal projection onto $L^2_0(-a,a)$, and $G_a=P_aGP_a$, Suzuki
defines $B_a=D^*G_aD$ on $\mathcal D(B_a)=H^1_0(-a,a)$, where
$D=i\,d/dx$ with Dirichlet boundary conditions, and proves that the
self-adjoint operator $A_a$ representing the localized Weil form
$Q^a_W$ is the Friedrichs extension of $B_a$
\cite[Thm.~1.1]{Suz26}; a first numerical study of the Hilbert space
derived from the Weil distribution appears in \cite{Suz26b}.  For
$v\in H^1_0(-a,a)$,
\begin{equation}\label{eq:form}
Q^a_W(v)=\langle B_av,v\rangle_{L^2}
=\int_{-a}^{a}\!\!\int_{-a}^{a} g(x-y)\,v'(x)\,v'(y)\,dx\,dy .
\end{equation}

\subsection{The TFPT window form}
The TFPT prime-front program is the part of the TFPT verification
suite that measures positivity surfaces of the Weil explicit formula
on declared finite windows; every step below is frozen in the
verification module \texttt{v563\_paper2\_readouts.py}.  It fixes a
prime-power zone $n_0$, sets $\alpha=\log n_0$, chooses a lattice of
$M$ cells of width $D=2\alpha/M$ (``frame~A, $\nu=4$'' names the
suite's declared window convention; a reader outside the suite may
take the explicit values of $\alpha,M,D$ quoted below as given), and
assembles two lag vectors from tent functions
$\tau_d(t)=(1-|t-dD|/D)_+$:
\begin{align}
c_{\mathrm{ar}}[d]&=-\int_0^\infty \tau_d(t)\,\rho(t)\,dt
\quad(d\ge1),\qquad
\rho(t)=\frac{e^{-t/2}}{1-e^{-2t}},\label{eq:car}\\
c_{\mathrm{at}}[d]&=-\sum_{n\le e^{2\alpha}}\frac{\Lam(n)}{\sqrt n}
\bigl(\tau_d(\log n)+\tau_d(-\log n)\bigr),\label{eq:cat}
\end{align}
with the $d=0$ cell of \eqref{eq:car} regularized by the declared
near-cell scheme with $\delta_0$-weight $-(\gamma+\log\pi)$ and
regularizer $e^{-3w/2}$.  The window matrix is the odd (parity)
compression of the Toeplitz form of $c=c_{\mathrm{ar}}+c_{\mathrm{at}}$,
\begin{equation}\label{eq:window}
(A_h)_{rs}=c_{|r-s|}-c_{(M-1)-r-s},\qquad r,s=0,\dots,h-1,\quad h=M/2,
\end{equation}
and the measured surface is $Q_T(v)=v^{\mathsf T}A_hv$.  On the
declared derivation window, $h=184$, $M=368$,
$\alpha=\log 16=2.772588722240$, $D=0.015068416969$, with $70$ atoms
$\log n\le2\alpha$.

Reading a coefficient vector $V\in\mathbb R^M$ as the hat expansion
$\Phi_V=\sum_j V_j\varphi_j$ on the half-integer lattice
$p_j=(j-\tfrac{M-1}2)D$, the odd sector $V_{M-1-j}=-V_j$ of
\eqref{eq:window} is exactly the space of odd piecewise-linear
functions in $H^1_0(-\tilde a,\tilde a)$, $\tilde a=\alpha+D/2$.

\section{The identification}

\subsection{The projection bookkeeping closes trivially}

\begin{lemma}\label{lem:proj}
(i) For every $v\in H^1_0(-a,a)$ one has $\int_{-a}^a v'=v(a)-v(-a)=0$;
hence $u=Dv\in L^2_0(-a,a)$ automatically, $P_a(Dv)=Dv$, and
$\langle G_aDv,Dw\rangle=\langle GDv,Dw\rangle$ for all
$v,w\in H^1_0(-a,a)$.  In particular Suzuki's $L^2_0$ constraint
imposes \emph{no} linear condition on the coefficient vectors of any
$H^1_0$-conforming Galerkin family, and the Krein-kernel replacement
$g(x-y)\mapsto g(x-y)-g(x)-g(-y)+g(0)$ leaves the Galerkin matrix
unchanged.  (ii) The TFPT parity sector consists of odd functions
$\Phi_V\in H^1_0(-\tilde a,\tilde a)$ with $\int\Phi_V=0$: an
\emph{extra} mean-zero condition on the test function itself, selecting
a parity subspace of the form domain (compatible with
$\lambda_a=\min(\lambda_a^+,\lambda_a^-)$,
\cite[Sec.~4.5]{Suz26}).
\end{lemma}

The point of (i) is that the mean-zero condition lives on the
\emph{differentiated} objects ($D:H^1_0\to L^2_0$ is a bijection,
\cite[Sec.~8.2]{Suz26}), where it holds identically; the point of (ii)
is that oddness is a $v$-side condition Suzuki does not need but which
does no harm.  Both statements are verified computationally in
\texttt{w1\_theorem\_probe.py} (P1.1--P1.3): the Krein correction
evaluates to $-2(\int gu)(\int u)+g(0)(\int u)^2$ with
$\int u=0$ exactly on the sector (must-break control: for the
non-mean-zero $u\equiv1$ the two kernels differ by $2.8357$).

\subsection{Three identities and the theorem}

Write $g''$ for the distributional second derivative and define the
\emph{Weil measure}
\begin{equation}\label{eq:weilmeasure}
\mathcal W:=g''+\bigl(e^{t/2}+e^{-t/2}\bigr)
=\underbrace{\rho(t)\,dt\ (t\neq0)\ +\ \text{origin
block}}_{\mathcal W_{\mathrm{sm}}}
\;+\;\underbrace{\sum_{n\ge2}\frac{\Lam(n)}{\sqrt n}
\bigl(\delta_{\log n}+\delta_{-\log n}\bigr)}_{\mathcal W_{\mathrm{at}}},
\end{equation}
where the origin block is the $\Pf$/$\delta_0$ content of the
$\tfrac12|t|\log|t|+A|t|$ head of \eqref{eq:origin} (total $\delta_0$
mass $2A+1=\log2\pi+\gamma$ in the $\Pf$ scheme,
\cite[Sec.~2.5]{Suz26}).

\begin{theorem}[W1, measure level]\label{thm:w1}
Let $g$ be \eqref{eq:screw} and $\mathcal W$ as in
\eqref{eq:weilmeasure}.  Then, on the declared window
($\alpha=\log16$, $M=368$) and portably in $D$ (Sec.~\ref{sec:port}):
\begin{enumerate}
\item[(i)] $g''=-(e^{t/2}+e^{-t/2})+\mathcal W$ as tempered
distributions; the pole block $-(e^{t/2}+e^{-t/2})$ (the $\zeta$-poles
$s=0,1$) is exactly finite-rank after parity compression (measured
rank~$1$).
\item[(ii)] \textbf{Tent reads.}  The TFPT arch lags are the point-route
tent reads of $-\mathcal W_{\mathrm{sm}}$ at \emph{every} lag, with no
constant and no scalar:
$c_{\mathrm{ar}}[d]=-\langle\mathcal W_{\mathrm{sm}},\tau_d\rangle$
(verified at $d=1,2,3,10,50$ to relative $3.4\times10^{-52}$); the
candidate origin constant
$\kappa:=c_{\mathrm{ar}}[0]+\tfrac1D\bigl[g_{\mathrm{sm}}(-D)
-2g_{\mathrm{sm}}(0)+g_{\mathrm{sm}}(D)\bigr]$ \emph{vanishes}
($|\kappa|<10^{-53}$; equivalently the closed identity
$\tfrac2D\int_0^Dt\rho\,dt+2\int_D^\infty\!\rho\,dt
+\tfrac{f(D)-\Phi(1,2,1/4)}{2D}=0$ with
$f(t)=e^{-t/2}\Phi(e^{-2t},2,\tfrac14)$): the v563 near-cell scheme
\emph{is} Suzuki's origin bookkeeping.  The atom lags \eqref{eq:cat}
are the literal tent reads of $-\mathcal W_{\mathrm{at}}$ (exact).
\item[(iii)] \textbf{B-spline reads.}  The hat-Galerkin lags of
$B_a=D^*G_aD$, i.e.
$c_{\mathrm{gal}}(|i-j|)=\iint g(x-y)\varphi_i'(x)\varphi_j'(y)$,
satisfy $c_{\mathrm{gal}}(d)=\mathrm{pole}_K(d)
-\langle\mathcal W,K_d\rangle$, where $K$ is the hat autocorrelation
(cubic B-spline with knot weights $(1,4,1)/6$, mass $D^2$) and
\[
\mathrm{pole}_K(d)=\frac{32\cosh(dD/2)\,(e^{D/2}+e^{-D/2}-2)^2}{D^2}.
\]
The smooth density read $-\langle\rho,K_d\rangle$ is exact for
$d\ge3$ (verified to $3.7\times10^{-31}$ at $d=5$); the atom layer has
the closed B-spline form
$-\sum_n\frac{\Lam(n)}{\sqrt n}[K(u_n-dD)+K(u_n+dD)]$ (verified against
kink-split double integrals to $4.2\times10^{-17}$); the boundary cells
carry three exact computable constants
$\tilde B(0)=1.295327231917$, $\tilde B(1)=-0.671894595706$,
$\tilde B(2)=0.017023960540$.
\item[(iv)] \textbf{One-scalar dictionary.}  Consequently
\[
c_{\mathrm{TFPT}}[d]
=\frac1D\Bigl(c_{\mathrm{gal}}(d)-\mathrm{pole}_K(d)\Bigr)
\]
up to (a) the derived second-order law
$R(d)=1+\tfrac{D^2}{12}\tfrac{\rho''}{\rho}(dD)+O(D^4)\to1+\tfrac1{6d^2}$
(window-independent) on the smooth layer, (b) the measure-level
re-binnings (atoms: tent vs.\ B-spline reads of the same data; smooth
near cells: exact read ratios of the same density), and (c) the exact
constants of (ii)--(iii).  On the common odd sector the two quadratic
forms agree on ten $16$-mode test vectors to per-vector relative
deviation $\le1.28\times10^{-10}$ (median $3.0\times10^{-11}$;
$16$-mode block operator-norm ratio $1.34\times10^{-11}$), and the
identification is \emph{sign-compatible with positivity}: the measured
positive-definiteness of $A_h$ and $Q^a_W\ge0$ sit on the same side.
\end{enumerate}
\end{theorem}

\emph{Proof of the smooth structure in (i) (five lines).}  Term-wise,
$e^{-t/2}\Phi(e^{-2t},2,\tfrac14)=\sum_{n\ge0}(n+\tfrac14)^{-2}
e^{-(2n+\frac12)t}$ and
$\frac{d^2}{dt^2}e^{-(2n+\frac12)t}
=(2n+\tfrac12)^2e^{-(2n+\frac12)t}
=4(n+\tfrac14)^2e^{-(2n+\frac12)t}$, so
\begin{equation}\label{eq:collapse}
\frac{d^2}{dt^2}\Bigl[e^{-t/2}\Phi(e^{-2t},2,\tfrac14)\Bigr]
=4\sum_{n\ge0}e^{-(2n+\frac12)t}=\frac{4e^{-t/2}}{1-e^{-2t}}=4\rho(t)
\qquad(t>0):
\end{equation}
the Hurwitz--Lerch block is a geometric series in disguise.  With the
prefactor $\tfrac14$ of \eqref{eq:screw} and
$\bigl[-4(e^{t/2}+e^{-t/2}-2)\bigr]''=-(e^{t/2}+e^{-t/2})$, off the
atoms $g''(t)=-(e^{t/2}+e^{-t/2})+\rho(t)$. \hfill$\square$

The $\delta_0$ closure in (ii) rests on the duplication identity
$\psi(\tfrac14)=-\gamma-3\log2-\pi/2$ and the scheme conversion
$k=\int_0^\infty\rho(w)(1-e^{-3w/2})\,dw
=\tfrac12(\psi(1)-\psi(\tfrac14))=\tfrac{3\log2+\pi/2}2$, whence
$-(\gamma+\log\pi)-2k=\psi(\tfrac14)-\log\pi$: the two $\delta_0$
conventions are one object in one scheme
(\texttt{w1\_boundary\_closure\_probe.py}, B1--B2, symbolic + $30$-digit
certificates).  All quantitative claims of
Theorem~\ref{thm:w1} are checks in \texttt{w1\_theorem\_probe.py}
(C0, P1--P3; $11/11$ pass, verdict \textsc{w1-theorem}).

\begin{remark}[Convention lock and verbatim transfer]
\label{rem:convention}
The reading of the final bracket of \eqref{eq:screw} is fixed by
Suzuki's own Section~2.2: the printed constant
$A=\tfrac12(\log2\pi-\psi(2))=0.707546\ldots$ of \eqref{eq:origin}
holds \emph{only} if the Lerch term carries the coefficient $+\tfrac14$
(machine check: with coefficient $+\tfrac14$ the origin read
reproduces $A$ to $1.2\times10^{-9}$ at $t=10^{-8}$; with coefficient
$-1$ it is off by $\approx56$ and $g(0)=-\tfrac54\Phi(1,2,\tfrac14)
=-21.4967\neq0$).  An earlier tranche of machine verifications (probe
files \texttt{w1\_dictionary\_probe.py},
\texttt{w1\_boundary\_closure\_probe.py},
\texttt{w1\_frozen\_windows\_probe.py},
\texttt{w1\_matrix\_identity\_probe.py}) used the kernel
$\tilde g:=g-\tfrac54e^{-|t|/2}\Phi(e^{-2|t|},2,\tfrac14)$ (Lerch
coefficient $-1$).  Since
$\tilde g_{\mathrm{sm}}''=-4\rho$ while $g_{\mathrm{sm}}''=+\rho$, one
has $c_{\mathrm{gal}}^{\mathrm{sm}}(\tilde g)
=-4\,c_{\mathrm{gal}}^{\mathrm{sm}}(g)$ \emph{exactly} (checked to
$3.7\times10^{-12}$), so every measured ratio of that tranche transfers
verbatim:
$c_{\mathrm{gal}}(\tilde g)/(-4Dc_{\mathrm{ar}})
\equiv c_{\mathrm{gal}}(g)/(+Dc_{\mathrm{ar}})$.  All numbers quoted
below from those probes are therefore valid for the true dictionary;
only the labels change ($-4D$ two-layer $\mapsto$ $+D$ one-scalar, and
the sign of the identification becomes positivity-compatible).
\end{remark}

\begin{remark}[Suite provenance]\label{rem:suite}
The probe files quoted in this note have been promoted to the
permanent, machine-checked verification suite of the TFPT repository
\cite{TFPT} as the modules \texttt{v630} (first contact: the atom
layer is literally Suzuki's prime measure), \texttt{v631} (per-lag
dictionary), \texttt{v640} (boundary closure), \texttt{v641}
(frozen-window portability), \texttt{v642} (matrix identity),
\texttt{v643} (the measure-level theorem of Sec.~2, including the
convention lock of Remark~\ref{rem:convention} as check C0.1), and
\texttt{v644} (the W2 slices of Sec.~\ref{sec:scope}).  The
\texttt{v631}/\texttt{v640}--\texttt{v642} chain was written in the
$\tilde g$-normalization of Remark~\ref{rem:convention}; the
convention was corrected the same day, see the dated erratum in the
suite (erratum blocks in \texttt{v631} and
\texttt{v640}--\texttt{v642}, machine evidence \texttt{v643} check
C0.1).  By the exact transfer identity above, every measured number of
that chain remains valid verbatim under the corrected convention.
\end{remark}

\subsection{The second-order law}
The tent and B-spline reads of a smooth density $f$ differ by derived
moment brackets (all coefficients symbolic):
$\int\tau_0t^{2k}=D,\ \tfrac{D^3}6,\ \tfrac{D^5}{15},\ \tfrac{D^7}{28},
\ \tfrac{D^9}{45}$ and
$\int K\,t^{2k}=D^2,\ \tfrac{D^4}3,\ \tfrac{3D^6}{10},\
\tfrac{17D^8}{42},\ \tfrac{31D^{10}}{45}$ for $2k=0,\dots,8$, giving
\begin{equation}\label{eq:momentlaw}
R(d):=\frac{\langle\rho,K_d\rangle}{D\,\langle\rho,\tau_d\rangle}
=1+\frac{D^2}{12}\,\frac{\rho''}{\rho}(dD)+O(D^4)
\;\longrightarrow\;1+\frac1{6d^2}\qquad(\text{window-independent}).
\end{equation}
On the declared window the exact ratios at $d=5,9,16$ are
$1.006766,\,1.001959,\,1.000585$; the $D^6$-truncated bracket explains
the excess $R-1$ to $3\times10^{-4}$, and the derived $D^8$ order
accounts for $77\%$--$99\%$ of the residual at $d=3,\dots,12$
(\texttt{w1\_theorem\_probe.py} P2.5): the remainder of the dictionary
is typed bracket truncation, not an unexplained conversion.

\section{Numerical verification}\label{sec:port}

\subsection{Portability across windows}
The dictionary was derived on $h=184$ and then \emph{frozen} (all
formulas closed in $D$, no knob) and executed unchanged on three fresh
frame-A windows (\texttt{w1\_frozen\_windows\_probe.py}, preregistered
kill test, $10/10$ pass, verdict \textsc{w1-dictionary-portable};
ratios transfer to the true dictionary by
Remark~\ref{rem:convention}):

\begin{center}
\small
\setlength{\tabcolsep}{4.5pt}
\begin{tabular}{lcccccccc}
\toprule
window & $M$ & $D$ & \multicolumn{4}{c}{smooth ratio
$c^{\mathrm{sm}}_{\mathrm{gal}}/(Dc_{\mathrm{ar}})$ at
$d=3,5,9,16$} & $|R(16)-1|$ & atoms\\
\midrule
$h=285$ & $570$ & $0.01731394$ & $1.0211$ & $1.0067$ & $1.0019$
& $1.0006$ & $5.8\times10^{-4}$ & $1.000000\,D^2$\\
$h=540$ & $1080$ & $0.00936342$ & $1.0213$ & $1.0069$ & $1.0020$
& $1.0006$ & $6.1\times10^{-4}$ & $1.000000\,D^2$\\
$h=997$ & $1994$ & $0.00603181$ & $1.0214$ & $1.0069$ & $1.0020$
& $1.0006$ & $6.2\times10^{-4}$ & $1.000000\,D^2$\\
\bottomrule
\end{tabular}
\end{center}

\noindent Here ``atoms'' is the per-atom Galerkin mass over the TFPT
tent mass ($\log2$ and $\log3$ agree to $<10^{-14}$): the atomic
re-binning constant is exactly the B-spline/tent mass ratio $D^2$, i.e.
the same $+1/D$ per-lag scalar as the smooth layer.  The moment
prediction \eqref{eq:momentlaw} explains the excess to $0.6\%$ on all
three windows; only $D$ changes.

\subsection{Form equality on the common odd sector}
On $h=184$, with ten odd $16$-mode test vectors (fixed seed) and the
dictionary applied in stages (\texttt{w1\_theorem\_probe.py} P3.1),
the $16$-mode block operator-norm ratios
$\lVert T(A_T-A_S)T'\rVert/\lVert TA_TT'\rVert$ are
\begin{center}
\begin{tabular}{lc}
\toprule
stage & norm ratio\\
\midrule
one-scalar $1/D$ & $4.37\times10^{-3}$\\
$+\ \tilde B(0..2)$ (exact boundary) & $3.00\times10^{-3}$\\
$+$ moment law \eqref{eq:momentlaw} & $2.72\times10^{-3}$\\
$+$ measure-level atom re-binning & $4.62\times10^{-6}$\\
$+$ measure-level smooth re-binning (near cells) & $1.34\times10^{-11}$\\
\bottomrule
\end{tabular}
\end{center}
and the per-vector deviations $|Q_T-Q_S|/|Q_T|$ reach at most
$1.28\times10^{-10}$ (median $3.0\times10^{-11}$).  With the derived
moment law alone (no exact near-cell ratios) the worst vector sits at
$1.02\times10^{-4}$, entirely the typed $D^8$ truncation at $d=3,4$
hitting near-cancellation denominators.  Nothing is fitted at any
stage.

\section{Scope: what is and is not claimed}\label{sec:scope}

This note establishes an \emph{identification of finite-dimensional
quadratic forms} and its exact measure-level skeleton.  It does
\emph{not} establish positivity of anything, and it makes no statement
about the Riemann Hypothesis.  The honest implication map is:

\begin{center}
\small
\begin{tabular}{@{}p{0.045\textwidth}p{0.33\textwidth}p{0.53\textwidth}@{}}
\toprule
step & statement & status\\
\midrule
W1 & window form $=$ Galerkin read of $B_a$ (explicit dictionary)
& \textbf{closed} at exact/certified level (this note;
\texttt{w1\_theorem\_probe.py}, $11/11$).\\
W2 & nested window spaces are form-dense in the odd form sector of
$A_a$; window spectra converge & \textbf{started}
(\texttt{w2\_form\_density\_probe.py}): FEM density with classical
rates verified ($H^1$ rate $1.000$, $L^2$ rate $2.000$ over
$M=92$--$736$); Rayleigh--Ritz eigenvalue convergence and
$\lambda_a$-profile measured; Mosco/norm-resolvent convergence
\emph{open}.\\
W3 & control in $a$: continuity/limits of $\lambda_a$, the
$a\to\infty$ strong-resolvent limit of \cite[Sec.~7]{Suz26}
& \emph{open}; \cite[Thm.~1.3]{Suz26} gives continuity of
$\lambda_a$ (no monotonicity claimed there).\\
W4 & positivity transfer: window positivity measurements
$\Rightarrow$ Weil positivity & \emph{not claimed}; would require
W2 $+$ W3 $+$ uniformity in the window family.  Weil positivity for
all $a$ is equivalent to RH \cite{Weil52,Yoshida92}; nothing here
touches that equivalence.\\
\bottomrule
\end{tabular}
\end{center}

The W2 slices measured in \texttt{w2\_form\_density\_probe.py} (suite
module \texttt{v644}, verdict \textsc{w2-slices-pass}, $7/7$) are,
beyond the FEM rates of the table: Rayleigh--Ritz on the nested dyadic
refinements $M=92/184/368/736$ at fixed $a=\log16$ gives lowest-three
eigenvalues decreasing monotonically from above in all six sequences
(full and odd sector), with Cauchy gaps shrinking at every step and
the parity identity $\lambda_{\mathrm{full}}
=\min(\lambda_{\mathrm{even}},\lambda_{\mathrm{odd}})$ holding to
$10^{-15}$ per refinement.  The $\lambda_a$ scan across five frame-A
windows $a=2.77,\dots,6.01$ at fixed $M=368$ gives $\lambda_a$ between
$+10^{-12}$ and $+2\times10^{-10}$ with per-window discretization gaps
$\sim10^{-9}$: the honest reading is $\lambda_a=0^+$ within
$\sim10^{-9}$ --- consistent with the continuity statement of
\cite[Thm.~1.3]{Suz26} and \emph{not} a sign statement --- with
ground-state parity even at every measured $a$ (reported;
\cite[Thm.~1.4]{Suz26} covers only small $a$).

Within W1 itself the remaining declared freedoms are exactly: the
finite-rank pole block (rank $1$ after parity compression), the typed
bracket truncation of \eqref{eq:momentlaw} beyond $D^8$, and the
floating-point/quadrature certification level stated with each check.

\begin{thebibliography}{9}
\bibitem{Suz26} M.~Suzuki,
\emph{Weil's quadratic form via the screw function},
arXiv:2606.09096 (2026).
\bibitem{Suz26b} M.~Suzuki,
\emph{On the Hilbert space derived from the Weil distribution},
Canadian J.\ Math.\ (to appear), arXiv:2607.24830.
\bibitem{Weil52} A.~Weil,
\emph{Sur les ``formules explicites'' de la th\'eorie des nombres
premiers}, Comm.\ S\'em.\ Math.\ Univ.\ Lund (1952), Tome
Suppl\'ementaire, 252--265.
\bibitem{Yoshida92} H.~Yoshida,
\emph{On Hermitian forms attached to zeta functions}, in: Zeta
functions in geometry (Tokyo, 1990), Adv.\ Stud.\ Pure Math.\ 21,
Kinokuniya, Tokyo (1992), 281--325.
\bibitem{TFPT} TFPT repository
(\url{https://fixpoint-theory.com}): verification modules
\texttt{v563\_paper2\_readouts.py}, \texttt{v630},
\texttt{v631}, \texttt{v640}--\texttt{v644}, and discovery probes
\texttt{w1\_dictionary\_probe.py},
\texttt{w1\_boundary\_closure\_probe.py},
\texttt{w1\_frozen\_windows\_probe.py},
\texttt{w1\_matrix\_identity\_probe.py},
\texttt{w1\_theorem\_probe.py},
\texttt{w2\_form\_density\_probe.py} (2026).
\end{thebibliography}

\end{document}
